Los problemas de ruteo de vehículos ocupan un lugar central en la investigación de
operaciones, ya que articulan decisiones de asignación, secuenciamiento y transporte
bajo restricciones operacionales. Entre sus variantes, el \textit{Capacitated Vehicle
Routing Problem} (CVRP) destaca como un modelo base para estudiar la distribución
física de mercancías desde un depósito central hacia clientes con demanda conocida,
cuando la capacidad de carga de cada vehículo es limitada
\parencite{toth2002vehicle}.

En el Informe~Nro.~1 este problema se abordó mediante métodos exactos de programación
matemática, capaces de certificar optimalidad sobre instancias pequeñas y medianas,
pero cuyo esfuerzo computacional crece de forma abrupta con el tamaño de la instancia.
El presente informe complementa ese trabajo abordando el CVRP mediante una
\textit{metaheurística}, \textit{Particle Swarm Optimization} (PSO), y contrastando
ambos enfoques sobre un conjunto común de instancias, según lo requerido por el
enunciado del trabajo de investigación.

\subsection{Contexto y motivación}

El interés práctico por el CVRP proviene de la necesidad de diseñar redes de
distribución con bajo costo total y alto nivel de servicio. En operaciones logísticas
reales, una planificación de rutas deficiente incrementa kilómetros recorridos,
horas-hombre, consumo de combustible y uso de flota, deteriorando la competitividad del
sistema. Desde una perspectiva de modelación, el CVRP representa la interacción entre
decisiones de partición de clientes en rutas y decisiones de ordenamiento de visitas
dentro de cada recorrido, lo que lo convierte en un problema estructuralmente complejo.

Dos contextos de aplicación ilustran esta relevancia. En distribución de retail y
comercio mayorista, un centro de distribución debe abastecer múltiples locales o puntos
de venta con una flota propia o subcontratada; aquí, el cumplimiento de capacidad y la
consolidación eficiente de carga son determinantes para reducir costos de transporte y
frecuencia de despachos. En recolección de residuos domiciliarios, los camiones deben
visitar sectores urbanos completos, respetando la capacidad de compactación antes de
retornar al depósito o estación de descarga; en este caso, la factibilidad operacional
de cada ruta depende directamente de la demanda acumulada por sector. En ambos
escenarios, el CVRP entrega una representación matemática suficientemente rica para
apoyar decisiones tácticas y operacionales.

\subsection{Fundamentos teóricos}

El CVRP pertenece a la familia de problemas $\mathcal{NP}$-hard y generaliza, entre
otros, al problema del viajante de comercio cuando la capacidad es suficientemente
grande o cuando se utiliza un solo vehículo \parencite{lenstra1981complexity}. Esta
dificultad combinatoria motiva, más allá de los métodos exactos, el uso de
\textit{metaheurísticas}: procedimientos de búsqueda de alto nivel que exploran el
espacio de soluciones guiados por reglas heurísticas, sacrificando la garantía de
optimalidad a cambio de obtener soluciones de buena calidad en tiempos acotados sobre
instancias que quedan fuera del alcance de los métodos exactos.

Las metaheurísticas se dividen, a grandes rasgos, en aquellas de \textit{trayectoria}
(refinan una única solución candidata, como la búsqueda tabú o el temple simulado) y
aquellas \textit{basadas en población}, que evolucionan un conjunto de soluciones
candidatas en paralelo. \textit{Particle Swarm Optimization} (PSO), introducido por
\textcite{kennedy1995particle}, pertenece a esta segunda familia: un \textit{enjambre}
de partículas explora el espacio de soluciones imitando el comportamiento colectivo de
organismos sociales (bandadas de aves, cardúmenes). Cada partícula $i$ mantiene una
posición $\mathbf{x}_i$, que codifica una solución candidata, y una velocidad
$\mathbf{v}_i$, que representa su desplazamiento en cada iteración; la velocidad se
actualiza combinando tres influencias: la inercia de su movimiento anterior, la
atracción hacia su mejor posición histórica (\textit{pbest}) y la atracción hacia la
mejor posición encontrada por el enjambre (\textit{gbest}), ponderadas respectivamente
por un peso de inercia y por coeficientes de aceleración cognitivo y social
\parencite{poli2007particle}.

PSO fue formulado originalmente para dominios \textit{continuos}: la posición de una
partícula es un vector real y su actualización es una simple suma vectorial
posición-más-velocidad. El CVRP, en cambio, es un problema \textit{combinatorio}: una
solución es un conjunto de rutas, no un punto en $\mathbb{R}^d$, y no existe una noción
directa de ``sumar una velocidad'' a una partición de clientes en rutas. Este desajuste
de dominios es el desafío central que debe resolverse para aplicar PSO al CVRP, y exige
definir un esquema de \textit{codificación} (cómo un vector real representa una
solución del CVRP) y su correspondiente procedimiento de \textit{decodificación} (cómo
recuperar un conjunto de rutas factible a partir de ese vector). La decisión de diseño
concreta adoptada en este trabajo para resolver ese desajuste se presenta en el
Capítulo~\ref{cap:metodo}, tras revisar en el Capítulo~\ref{cap:estado-arte} las
alternativas consideradas en la literatura.

\subsection{Objetivos}

El objetivo general de este informe es resolver el CVRP mediante la metaheurística
\textit{Particle Swarm Optimization} y comparar su desempeño con el método exacto
implementado en el Informe~Nro.~1, sobre instancias estándar de la literatura.

Los objetivos específicos son los siguientes:
\begin{enumerate}
    \item Revisar el estado del arte de las metaheurísticas aplicadas al CVRP,
          posicionando PSO dentro de dicho panorama.
    \item Definir un esquema de representación (codificación y decodificación) que
          permita aplicar PSO, algoritmo de naturaleza continua, sobre el espacio de
          soluciones combinatorio del CVRP.
    \item Implementar computacionalmente el algoritmo PSO para el CVRP, incluyendo el
          manejo de las restricciones de capacidad y la construcción de rutas
          factibles.
    \item Evaluar el método sobre instancias de los Sets E y A de CVRPLIB (las mismas
          empleadas en el Informe~1) y, adicionalmente, sobre instancias de mayor
          tamaño que exceden el alcance del método exacto.
    \item Contrastar los resultados de PSO frente al método exacto en términos de
          calidad de solución (gap respecto del óptimo o de la mejor solución
          conocida), tiempo de ejecución y escalabilidad.
\end{enumerate}

\subsection{Organización del informe}

El Capítulo~2 presenta la definición del problema y su formulación matemática sobre un
grafo no dirigido. El Capítulo~3 resume el estado del arte de las metaheurísticas para el
CVRP, posicionando PSO dentro de ese panorama. El Capítulo~4 describe el diseño del
algoritmo PSO, con énfasis en el esquema de codificación/decodificación, la
parametrización y el protocolo experimental. El Capítulo~5 reporta los resultados de la
experimentación computacional. El Capítulo~6 discute los hallazgos y desarrolla la
comparación entre el enfoque exacto y el metaheurístico. Finalmente, el Capítulo~7
presenta las conclusiones.
